%% =====================================================================
%%  Abstract and 要旨
%%  A Master's thesis written in English must have BOTH an English abstract
%%  and a Japanese 要旨 (laboratory rule). Each should fit on one page.
%%  Cover the background, problem, objective, method, results and conclusion.
%%  Do not use figures, tables, equations or citations in the abstract.
%%  Lines such as "%% Background" below are comments that show the structure
%%  (they do not appear in the PDF). Do not write such labels in the text.
%%  The 論文要旨 (2-8 pages) required by the Graduate School is a separate
%%  document; see README.md.
%% =====================================================================
\begin{abstractEN}
\begin{writingnote}
Write the abstract last, after the rest of the thesis, and do not include anything that is not in the main text.
The objective and the main results are essential: keep the background and the problem to one or two sentences,
and spend most of the words on the method and the results.
State the results concretely: what you found and under which conditions, preferably with the key numbers.
Avoid vague phrases such as ``the effectiveness was confirmed''.
End with what the results imply, including the applicability and limitations of the method.
Define abbreviations at first use, and use the same terms and numbers as in the main text, especially in the conclusion.
Use the present tense for the problem and for what this thesis does (``This thesis proposes \ldots'')
and the past tense for what was done in the experiments.
\end{writingnote}

%% Background
In recent years, XX has become increasingly important in \ldots.
%% Problem
However, existing methods have the problem that \ldots.

%% Objective and method
This thesis proposes a method for the YY problem using XX in order to \ldots.
The proposed method \ldots.

%% Results and conclusion
Computational experiments show that \ldots{} under the conditions \ldots.
These results indicate that the method is applicable to \ldots, with the limitation that \ldots.
\end{abstractEN}

\begin{abstractJA}
\begin{writingnote}
Write the same content as the Abstract in Japanese, in the plain style (である調) and with the punctuation 「，」 and 「．」, as in the example below.
If you use machine translation, check that the technical terms and numbers match those in the Abstract and the main text.
\end{writingnote}

%% Abstract と同じ内容を日本語で書く．
近年，……の分野では，○○が重要になっている．
しかし，従来の方法には……という課題がある．
そこで本研究では，……を目的として，○○を用いた△△問題の解法を提案する．
数値実験により，……の条件で……という結果が得られた．
以上より，提案手法の適用条件として……，限界として……が明らかになった．
\end{abstractJA}
